\section{Introduction}
\label{sec:intro}

% Figure 1 --- the overview teaser placed right after the Introduction heading:
% the end-to-end DynamicMCPBench pipeline.
% AUTO-GENERATED by paper/regenerate.py -- do not hand-edit.
% id: fig:pipeline
% status: manual
% reason: manual block diagram -- author from method section. Source-of-truth: docs/CONCEPT.md.
\begin{figure}[t]
  \centering
  \fbox{\parbox{0.9\columnwidth}{\centering \textbf{fig:pipeline} -- placeholder.\\
  Status: manual.\\
  Gating step: ---}}
  \caption{DynamicMCPBench pipeline: server manifest $\to$ goal-gen $\to$ explorer $\to$ distiller $\to$ evaluator (Tier-1 / Tier-2) $\to$ report; the parallel `refresh` arrow for the living-bench loop.}
  \label{fig:pipeline}
\end{figure}


% TODO(intro-motivation): LLM agents over Model Context Protocol (MCP) servers
%   are a real deployment surface; their benchmarks are not. The two prevailing
%   "did the agent do the right thing?" proxies both wobble under live, stateful
%   data --- final-answer string matching and tool-list recall (the latter
%   diagnosed by prior art as ~50% noise).

% TODO(intro-pivot): change the organizing primitive from answer / tool set to
%   the *execution trace* --- the only object that is unambiguously observable
%   (we recorded it) and faithful (every tool in a trace was actually invoked).

% TODO(intro-preview): one paragraph previewing the forward-generate then
%   effect-score pipeline (Fig.~\ref{fig:pipeline}) and the headline ranking
%   inversion (\S\ref{sec:results}).

% TODO(intro-contributions): "Our contributions can be summarized as follows:"
%   1. A forward generator (explore -> distill) that compiles a live successful
%      trajectory into a path-agnostic TaskSpec with non-trivial equivalence sets.
%   2. A trace/effect alignment scorer (Tier-1 deterministic + Tier-2 judge) that
%      never grades the final answer.
%   3. A trace-property failure model isolating task dynamism as the dominant driver.
%   4. A living-bench refresh protocol over a reproducible, sandboxed MCP substrate.
